% GENERATED FILE. Edit canonical lessons, then run npm run generate:books.
% canonical-source-hash: fnv1a64:588f283bca3048b5
% canonical-lessons: MR-C501-sparsa, MR-C501-mithi-marne, MR-C501-gonjarne, MR-C501-thopatne, MR-C501-kurvalne

\chapter{Hugs and Pats}
\label{ch:mr-hugs-and-pats}
\hlchaptermodality{501}

\begin{chapteropening}
\textbf{By the end of this chapter:} \emph{I can talk about a touch, a hug, a pat and a caress.}

Complete the last lesson of chapter 501: I can talk about a touch, a hug, a pat and a caress.
\end{chapteropening}

\section[\texorpdfstring{\mr{स्पर्श} (sparśa)}{स्पर्श (sparśa)}]{\mr{स्पर्श} (sparśa) — a touch}

\label{lesson:MR-C501-sparsa}



\begin{quote}
Before the new one: say the Marathi for to pamper, then the Marathi for to adopt.
\end{quote}

\subsection*{You'll want to know: \mr{स्पर्श}}
\textbf{\mr{स्पर्श}} — \emph{sparśa} — \textquotedblleft{}a touch\textquotedblright{}.

\begin{quote}
\mr{आईच्या} \mr{स्पर्शाने} \mr{बाळ} \mr{शांत} \mr{झालं}.
\end{quote}

\begin{quote}
\emph{āīcyā sparśāne bāḷ śānt jhālaṁ.}
\end{quote}

\begin{quote}
The baby calmed down at its mother's touch.
\end{quote}

Masculine.

\mr{स्पर्श} \mr{करणे}, to touch.

\begin{cousinweb}
\textbf{\mr{स्पर्श}} (\emph{sparśa}) is a tatsama, taken from Sanskrit \emph{sparśa}, \textquotedblleft{}touch\textquotedblright{} (CDIAL 13809).
\end{cousinweb}

\begin{grammarlens}[title={What you've built}]
One more word for describing what you lived and felt.
\end{grammarlens}

\subsection*{Your turn}
\begin{itemize}
\raggedright
  \item \emph{Say it:} \emph{sparśa}
  \item \emph{Say it:} \emph{sparśa}, once more
  \item \emph{Say it:} \emph{āīcyā sparśāne bāḷ śānt jhālaṁ.} — The baby calmed down at its mother's touch.
  \item \emph{Recall:} say the Marathi for to pamper, then the Marathi for to adopt, then say \emph{sparśa} again
\end{itemize}

\begin{tcolorbox}[breakable,colback=teal!4,colframe=teal!35!black,title={Before you move on}]
What does \mr{स्पर्श} mean? (\textquotedblleft{}A touch\textquotedblright{}.) Say it once more.
\end{tcolorbox}
\section[\texorpdfstring{\mr{मिठी} \mr{मारणे} (miṭhī mārṇe)}{मिठी मारणे (miṭhī mārṇe)}]{\mr{मिठी} \mr{मारणे} (miṭhī mārṇe) — to hug}

\label{lesson:MR-C501-mithi-marne}



\begin{quote}
Before the new one: say the Marathi for to adopt, then the Marathi for a touch.
\end{quote}

\subsection*{You'll want to know: \mr{मिठी} \mr{मारणे}}
\textbf{\mr{मिठी} \mr{मारणे}} — \emph{miṭhī mārṇe} — \textquotedblleft{}to hug\textquotedblright{}.

\begin{quote}
\mr{मित्राला} \mr{पाहून} \mr{त्याने} \mr{मिठी} \mr{मारली}.
\end{quote}

\begin{quote}
\emph{mitrālā pāhūn tyāne miṭhī mārlī.}
\end{quote}

\begin{quote}
When he saw his friend he gave him a hug.
\end{quote}

Past, with -\mr{ने} (he …): \mr{त्याने} \mr{मिठी} \mr{मारली}, \emph{tyāne miṭhī mārlī}.

\mr{मिठी}, a hug; the one hugged takes -\mr{ला}: \mr{आईला} \mr{मिठी} \mr{मारणे}.

\begin{cousinweb}
\textbf{\mr{मिठी} \mr{मारणे}} (\emph{miṭhī mārṇe}) is a Marathi phrase built on \mr{मिठी}, \textquotedblleft{}a hug\textquotedblright{}, whose origin is not settled, with \mr{मारणे} as the light verb.
\end{cousinweb}

\begin{grammarlens}[title={What you've built}]
One more word for describing what you lived and felt.
\end{grammarlens}

\subsection*{Your turn}
\begin{itemize}
\raggedright
  \item \emph{Say it:} \emph{miṭhī mārṇe}
  \item \emph{Say it:} \emph{miṭhī mārṇe}, once more
  \item \emph{Say it:} \emph{mitrālā pāhūn tyāne miṭhī mārlī.} — When he saw his friend he gave him a hug.
  \item \emph{Recall:} say the Marathi for to adopt, then the Marathi for a touch, then say \emph{miṭhī mārṇe} again
\end{itemize}

\begin{tcolorbox}[breakable,colback=teal!4,colframe=teal!35!black,title={Before you move on}]
What does \mr{मिठी} \mr{मारणे} mean? (\textquotedblleft{}To hug\textquotedblright{}.) Say it once more.
\end{tcolorbox}
\section[\texorpdfstring{\mr{गोंजारणे} (goñjārṇe)}{गोंजारणे (goñjārṇe)}]{\mr{गोंजारणे} (goñjārṇe) — to stroke; to pet}

\label{lesson:MR-C501-gonjarne}



\begin{quote}
Before the new one: say the Marathi for a touch, then the Marathi for to hug.
\end{quote}

\subsection*{You'll want to know: \mr{गोंजारणे}}
\textbf{\mr{गोंजारणे}} — \emph{goñjārṇe} — \textquotedblleft{}to stroke; to pet\textquotedblright{}.

\begin{quote}
\mr{मुलीने} \mr{मांजराला} \mr{गोंजारलं}.
\end{quote}

\begin{quote}
\emph{mulīne māñjrālā goñjārlaṁ.}
\end{quote}

\begin{quote}
The girl stroked the cat.
\end{quote}

Past, with -\mr{ने} (he …): \mr{त्याने} \mr{मांजराला} \mr{गोंजारलं}, \emph{tyāne māñjrālā goñjārlaṁ}.

Soft strokes for a pet or a child; the one stroked takes -\mr{ला}.

\begin{cousinweb}
\textbf{\mr{गोंजारणे}} (\emph{goñjārṇe}) is a Marathi verb whose origin is not settled.
\end{cousinweb}

\begin{grammarlens}[title={What you've built}]
One more word for describing what you lived and felt.
\end{grammarlens}

\subsection*{Your turn}
\begin{itemize}
\raggedright
  \item \emph{Say it:} \emph{goñjārṇe}
  \item \emph{Say it:} \emph{goñjārṇe}, once more
  \item \emph{Say it:} \emph{mulīne māñjrālā goñjārlaṁ.} — The girl stroked the cat.
  \item \emph{Recall:} say the Marathi for a touch, then the Marathi for to hug, then say \emph{goñjārṇe} again
\end{itemize}

\begin{tcolorbox}[breakable,colback=teal!4,colframe=teal!35!black,title={Before you move on}]
What does \mr{गोंजारणे} mean? (\textquotedblleft{}To stroke; to pet\textquotedblright{}.) Say it once more.
\end{tcolorbox}
\section[\texorpdfstring{\mr{थोपटणे} (thopaṭṇe)}{थोपटणे (thopaṭṇe)}]{\mr{थोपटणे} (thopaṭṇe) — to pat}

\label{lesson:MR-C501-thopatne}



\begin{quote}
Before the new one: say the Marathi for to hug, then the Marathi for to stroke.
\end{quote}

\subsection*{You'll want to know: \mr{थोपटणे}}
\textbf{\mr{थोपटणे}} — \emph{thopaṭṇe} — \textquotedblleft{}to pat\textquotedblright{}.

\begin{quote}
\mr{आई} \mr{बाळाला} \mr{थोपटून} \mr{झोपवते}.
\end{quote}

\begin{quote}
\emph{āī bāḷālā thopṭūn jhopavte.}
\end{quote}

\begin{quote}
Mother pats the baby to sleep.
\end{quote}

Past, with -\mr{ने} (he …): \mr{त्याने} \mr{बाळाला} \mr{थोपटलं}, \emph{tyāne bāḷālā thopaṭlaṁ}.

\mr{पाठ} \mr{थोपटणे}, to pat someone on the back.

\begin{cousinweb}
\textbf{\mr{थोपटणे}} (\emph{thopaṭṇe}) is usually traced to the expressive \emph{*thapp}, \textquotedblleft{}slap, pat\textquotedblright{}, under which CDIAL lists Marathi \emph{thābṇẽ}, \textquotedblleft{}to pat\textquotedblright{} (CDIAL 6091).
\end{cousinweb}

\begin{grammarlens}[title={What you've built}]
One more word for describing what you lived and felt.
\end{grammarlens}

\subsection*{Your turn}
\begin{itemize}
\raggedright
  \item \emph{Say it:} \emph{thopaṭṇe}
  \item \emph{Say it:} \emph{thopaṭṇe}, once more
  \item \emph{Say it:} \emph{āī bāḷālā thopṭūn jhopavte.} — Mother pats the baby to sleep.
  \item \emph{Recall:} say the Marathi for to hug, then the Marathi for to stroke, then say \emph{thopaṭṇe} again
\end{itemize}

\begin{tcolorbox}[breakable,colback=teal!4,colframe=teal!35!black,title={Before you move on}]
What does \mr{थोपटणे} mean? (\textquotedblleft{}To pat\textquotedblright{}.) Say it once more.
\end{tcolorbox}
\section[\texorpdfstring{\mr{कुरवाळणे} (kurvāḷṇe)}{कुरवाळणे (kurvāḷṇe)}]{\mr{कुरवाळणे} (kurvāḷṇe) — to caress}

\label{lesson:MR-C501-kurvalne}



\begin{quote}
Before the new one: say the Marathi for to stroke, then the Marathi for to pat.
\end{quote}

\subsection*{You'll want to know: \mr{कुरवाळणे}}
\textbf{\mr{कुरवाळणे}} — \emph{kurvāḷṇe} — \textquotedblleft{}to caress\textquotedblright{}.

\begin{quote}
\mr{आजीने} \mr{प्रेमाने} \mr{मला} \mr{कुरवाळलं}.
\end{quote}

\begin{quote}
\emph{ājīne premāne malā kurvāḷlaṁ.}
\end{quote}

\begin{quote}
Grandmother caressed me lovingly.
\end{quote}

Past, with -\mr{ने} (he …): \mr{त्याने} \mr{कुत्र्याला} \mr{कुरवाळलं}, \emph{tyāne kutryālā kurvāḷlaṁ}.

A loving stroke of the head or face; the one caressed takes -\mr{ला}.

\begin{cousinweb}
\textbf{\mr{कुरवाळणे}} (\emph{kurvāḷṇe}) is a Marathi verb whose origin is not settled.
\end{cousinweb}

\begin{grammarlens}[title={What you've built}]
One more word for describing what you lived and felt.
\end{grammarlens}

\subsection*{Your turn}
\begin{itemize}
\raggedright
  \item \emph{Say it:} \emph{kurvāḷṇe}
  \item \emph{Say it:} \emph{kurvāḷṇe}, once more
  \item \emph{Say it:} \emph{ājīne premāne malā kurvāḷlaṁ.} — Grandmother caressed me lovingly.
  \item \emph{Recall:} say the Marathi for to stroke, then the Marathi for to pat, then say \emph{kurvāḷṇe} again
\end{itemize}

\begin{tcolorbox}[breakable,colback=teal!4,colframe=teal!35!black,title={Before you move on}]
What does \mr{कुरवाळणे} mean? (\textquotedblleft{}To caress\textquotedblright{}.) Say it once more.
\end{tcolorbox}
